\subsection{Recorridos básicos: BFS y DFS}
\label{alg:2-9}

Casi todo problema de grafos empieza con uno de estos dos recorridos. Aquí van las
versiones mínimas, listas para copiar; las entradas de grafos \ref{alg:4-11} y \ref{alg:4-12} tienen las
completas (reconstrucción del camino, 0-1 BFS, DFS con preorden y postorden).

\subsubsection{Leer el grafo}

La representación estándar es una lista de adyacencia: \texttt{adj[u]} guarda los
vecinos de \texttt{u}.

\begin{lstlisting}[language=Python]
import sys
from collections import deque

data = sys.stdin.buffer.read().split()
n, m = int(data[0]), int(data[1])
adj = [[] for _ in range(n + 1)]          # nodos 1..n
for i in range(m):
    a, b = int(data[2 + 2*i]), int(data[3 + 2*i])
    adj[a].append(b)
    adj[b].append(a)                      # quitar si es dirigido
\end{lstlisting}

\vspace{0.3cm}
\subsubsection{BFS: distancias desde un origen}

Recorre por capas: la primera vez que llega a un nodo lo hace por el camino con
menos aristas. Sirve para caminos mínimos sin pesos, "mínimo número de pasos" y
todo lo que sea explorar por niveles. Costo \bigO{V + E}.

\begin{lstlisting}[language=Python]
def bfs(s):
    dist = [-1] * (n + 1)                 # -1 = no visitado
    dist[s] = 0
    q = deque([s])
    while q:
        u = q.popleft()
        for v in adj[u]:
            if dist[v] == -1:
                dist[v] = dist[u] + 1
                q.append(v)
    return dist
\end{lstlisting}

\graybold{Cuidado:} la cola debe ser un \texttt{deque}. Con una lista,
\texttt{q.pop(0)} cuesta \bigO{n} y vuelve cuadrático el recorrido.

\vspace{0.3cm}
\subsubsection{BFS en grilla}

La grilla es un grafo implícito: cada casilla libre conecta con sus cuatro vecinas.

\begin{lstlisting}[language=Python]
def bfs_grilla(g, sr, sc):                # g: lista de strings, '#' = pared
    R, C = len(g), len(g[0])
    dist = [[-1] * C for _ in range(R)]
    dist[sr][sc] = 0
    q = deque([(sr, sc)])
    while q:
        r, c = q.popleft()
        for dr, dc in ((1, 0), (-1, 0), (0, 1), (0, -1)):
            nr, nc = r + dr, c + dc
            if (0 <= nr < R and 0 <= nc < C and g[nr][nc] != '#'
                    and dist[nr][nc] == -1):
                dist[nr][nc] = dist[r][c] + 1
                q.append((nr, nc))
    return dist
\end{lstlisting}

La matriz se crea con \texttt{[[-1] * C for \_ in range(R)]} y no con
\texttt{[[-1] * C] * R}: esta segunda forma repite la MISMA fila $R$ veces, y
modificar una casilla las cambia todas.

\vspace{0.3cm}
\subsubsection{DFS iterativo: componentes conexas}

En Python conviene el DFS con pila explícita: no depende del límite de recursión ni
de la versión del intérprete.

\begin{lstlisting}[language=Python]
def componentes():
    comp = [0] * (n + 1)                  # 0 = sin visitar
    c = 0
    for s in range(1, n + 1):
        if comp[s]:
            continue
        c += 1
        comp[s] = c
        pila = [s]
        while pila:
            u = pila.pop()
            for v in adj[u]:
                if not comp[v]:
                    comp[v] = c
                    pila.append(v)
    return c, comp                        # cantidad y etiqueta de cada nodo
\end{lstlisting}

Marca cada nodo al meterlo en la pila, así ninguno entra dos veces. Alcanza para
componentes, alcanzabilidad y detección de bipartito; cuando hace falta el orden
exacto de entrada y salida de cada nodo (puentes, orden topológico por DFS), la
plantilla es la de la entrada \ref{alg:4-12}.

\vspace{0.3cm}
\subsubsection{DFS recursivo (solo con poca profundidad)}

Es la forma más corta de escribir, y basta cuando el grafo es pequeño:

\begin{lstlisting}[language=Python]
sys.setrecursionlimit(1 << 25)
vis = [False] * (n + 1)

def dfs(u):
    vis[u] = True
    for v in adj[u]:
        if not vis[v]:
            dfs(v)
\end{lstlisting}

Con miles de niveles de profundidad (un camino largo, por ejemplo), revisar la nota
sobre recursión del módulo \texttt{sys} (\ref{sub:sys}), o usar la versión
iterativa.
